\subsection{Beschreibung von idealen Eigenschaften}
% Alt: Aufgabe 8
\Aufgabe{
    \label{Aufg8}
    Beschreiben Sie folgende Kennzeichen eines idealen Operationsverstärkers:
    \begin{itemize}
        \item [ \bf a)] Eingangswiderstand
        \item [ \bf b)] Ausgangswiderstand
        \item [ \bf c)] Differenzverstärkung
        \item [ \bf d)] Eingangsoffsetspannung
    \end{itemize}
}

\Loesung{
    \begin{itemize}
        \item [ \bf a)] Eingangswiderstand: unendlich, um die Spannungsquelle wenig zu belasten
        \item [ \bf b)] Ausgangswiderstand: $0\,\Omega$, da OP so viel Leistung wie möglich abgeben soll
        \item [ \bf c)] Differenzverstärkung: unendlich, um kleine Signale zu verstärken
        \item [ \bf d)] Eingangsoffsetspannung: $0\,\mathrm{V}$, damit OP bei keinem Eingangssignal auch keine Ausgangsspannung ausgibt
    \end{itemize}
    Siehe: Tabelle \ref{tab:Ideale und typische Operationsverstärkereigenschaften}
}